% Placeholder for the practitioner article (SREcon extended abstract,
% USENIX ;login: article, or ACM Queue). Length target: 3-6 pages.

\documentclass[11pt]{article}
\usepackage[margin=1in]{geometry}
\usepackage{amsmath}
\usepackage{booktabs}
\usepackage{hyperref}

\title{Five Golden Signals for AI Workloads on GPU Clusters\\
(a practitioner article, draft)}
\author{Andrew Espira \and Sharath Kumar}
\date{Draft --- \today}

\begin{document}
\maketitle

\begin{abstract}
The four golden signals travel poorly to AI workloads on shared GPU
clusters. GPU utilization is a lying metric, HTTP errors miss silent
quality regressions, request latency misses queue wait and token tails,
and QPS is undefined for training. We propose five golden signals SREs
can start monitoring today, each with a copy-pasteable SLO template.
\end{abstract}

\section{The problem}
\emph{TODO. Two-paragraph intro. Golden signals worked for stateless
services; not for AI on GPUs.}

\section{Why classical four break --- one paragraph per signal}
\emph{TODO. See \texttt{docs/CRITIQUE.md}.}

\section{Five signals to monitor today}
\emph{TODO. Each signal capped at one page: definition,
what-to-alert-on, what-NOT-to-alert-on. See
\texttt{docs/SIGNALS\_TAXONOMY.md}.}

\section{SLO templates you can copy}
\emph{TODO. Prometheus / OpenTelemetry snippets from
\texttt{docs/SLO\_TEMPLATES.md}.}

\section{A short case study}
\emph{TODO. One failure mode from each of the four environments.}

\section{Open questions}
\emph{TODO. Honest list of what we don't know yet.}

\section*{Acknowledgments}
Framing conversations at Google SRE and Datadog. \emph{Names to be
added once the individuals involved confirm they are comfortable
being credited.}

\end{document}
